% Discussion
\section{Discussion}
\label{sec:discussion}

Locality weighting changes the measure by suppressing the contribution of
sectors far from the observable patch in field space.
This differs from global counting, where the most populous basin dominates
regardless of locality.

Cutoff stability matters because a measure with uncontrolled cutoff dependence
does not produce well-defined predictions.
The improvement of \ISR over global census is numerical and empirical, not a
theorem, but it demonstrates that locality weighting is a viable direction
for measure construction.

Basin boundaries limit locality continuity.
As the initial-condition spread $\sigma$ increases, more trajectories land in
distant basins, diluting the locality signal.
This is physical---there is no locality principle that guarantees nearby
field values share the same vacuum---and it motivates the diagnostics we
introduce.

The kernel-weighted sector averaging (Sec.~\ref{sec:results}) loosely echoes
the bookkeeping role of coarse-graining in effective field theory, but no
physical running-coupling interpretation is implied.

Several directions for future work are natural:
\begin{itemize}
  \item Replacing the ancestry proxy with a true causal ancestry graph
        based on branching history.
  \item Extending to multi-field landscapes with more realistic potentials.
  \item Incorporating Coleman--De Luccia transition rates between basins.
  \item Adding explicit decoherence or an open-system treatment of the
        environment.
  \item Bayesian model comparison against other measure proposals using
        synthetic observations.
  \item Analytic investigation of convergence conditions for different
        kernel families.
\end{itemize}
